\Askhsh{21193}{4}
\begin{ylh}{1}
    \item 11.5. Μήκος τόξου.
\end{ylh}
\begin{wrapfigure}[24]{r}{6cm}
    \centering
    \begin{tikzpicture}[scale=0.8]
        \def\R{0.6} % Ακτίνα

        % Ορισμός κέντρων κύκλων (κορυφές τριγώνου ΑΒΓ)
        \tkzDefPoint(0,3){A}
        \tkzDefPoint(5,3){Gamma}
        \tkzDefPoint(2.5,0){B}

        % Σχεδίαση τριγώνου που συνδέει τα κέντρα (πλευρές α, β, γ)
        \tkzDrawSegment[line width=0.5pt, color=secondarycolor](A,Gamma)
        \tkzDrawSegment[line width=0.5pt, color=secondarycolor](Gamma,B)
        \tkzDrawSegment[line width=0.5pt, color=secondarycolor](B,A)

        % Ετικέτες για τα μήκη των πλευρών
        \tkzLabelSegment[above, pos=0.5](A,Gamma){$\beta$}
        \tkzLabelSegment[right, pos=0.5](Gamma,B){$\alpha$}
        \tkzLabelSegment[left, pos=0.5](B,A){$\gamma$}

        % Σχεδίαση των κύκλων
        \tkzDrawCircle[R](A,\R cm)
        \tkzDrawCircle[R](Gamma,\R cm)
        \tkzDrawCircle[R](B,\R cm)

        % Ετικέτες για τα κέντρα
        \tkzDrawPoints(A,B,Gamma)
        \tkzLabelPoint[above left](A){A}
        \tkzLabelPoint[below](B){B}
        \tkzLabelPoint[above right](Gamma){$\Gamma$}

        % Ορισμός σημείων εφαπτόμενης για τα τμήματα του ιμάντα
        % Τμήμα ΛM (παράλληλο προς AΓ) - προς τα πάνω
        \tkzDefPointBy[translation=from A to Gamma](A) \tkzGetPoint{vAG}
        \tkzDefPointBy[rotation=center A angle 90](vAG) \tkzGetPoint{vAGperp} % Αντιωρολογιακή περιστροφή για "πάνω"
        \tkzDefPointWith[linear, K=\R](A,vAGperp) \tkzGetPoint{Lambda}
        \tkzDefPointWith[linear, K=\R](Gamma,vAGperp) \tkzGetPoint{Mu}

        % Τμήμα NP (παράλληλο προς ΓB) - προς τα κάτω-δεξιά
        \tkzDefPointBy[translation=from Gamma to B](Gamma) \tkzGetPoint{vGB}
        \tkzDefPointBy[rotation=center Gamma angle -90](vGB) \tkzGetPoint{vGBperp} % Ωρολογιακή περιστροφή για "κάτω-δεξιά"
        \tkzDefPointWith[linear, K=\R](Gamma,vGBperp) \tkzGetPoint{Nu}
        \tkzDefPointWith[linear, K=\R](B,vGBperp) \tkzGetPoint{Rho}

        % Τμήμα ΣK (παράλληλο προς BA) - προς τα κάτω-αριστερά
        \tkzDefPointBy[translation=from B to A](B) \tkzGetPoint{vBA}
        \tkzDefPointBy[rotation=center B angle -90](vBA) \tkzGetPoint{vBAperp} % Ωρολογιακή περιστροφή για "κάτω-αριστερά"
        \tkzDefPointWith[linear, K=\R](B,vBAperp) \tkzGetPoint{Sigma}
        \tkzDefPointWith[linear, K=\R](A,vBAperp) \tkzGetPoint{Kappa}

        % Σχεδίαση των ευθύγραμμων τμημάτων του ιμάντα
        \tkzDrawSegment[line width=2pt, color=maincolor](Lambda,Mu)
        \tkzDrawSegment[line width=2pt, color=maincolor](Nu,Rho)
        \tkzDrawSegment[line width=2pt, color=maincolor](Sigma,Kappa)

        % Σχεδίαση ακτίνων προς τα σημεία εφαπτόμενης (λεπτές γραμμές)
        \tkzDrawSegment[line width=0.5pt, color=secondarycolor](A,Lambda)
        \tkzDrawSegment[line width=0.5pt, color=secondarycolor](Gamma,Mu)
        \tkzDrawSegment[line width=0.5pt, color=secondarycolor](Gamma,Nu)
        \tkzDrawSegment[line width=0.5pt, color=secondarycolor](B,Rho)
        \tkzDrawSegment[line width=0.5pt, color=secondarycolor](B,Sigma)
        \tkzDrawSegment[line width=0.5pt, color=secondarycolor](A,Kappa)

        % Σχεδίαση και επισήμανση των σημείων εφαπτόμενης
        \tkzDrawPoints(Lambda,Mu,Nu,Rho,Sigma,Kappa)
        \tkzLabelPoint[above](Lambda){$\Lambda$}
        \tkzLabelPoint[above](Mu){M}
        \tkzLabelPoint[below right](Nu){N}
        \tkzLabelPoint[below](Rho){P}
        \tkzLabelPoint[below left](Sigma){$\Sigma$}
        \tkzLabelPoint[left](Kappa){K}

        % Σημείωση και επισήμανση των γωνιών ω, θ, φ (KÂΛ, ΣBΡ, MΓN)
        % Γωνία KÂΛ = ω (γύρω από το A)
        \tkzMarkAngle[fill=gray!30, opacity=0.8, size=0.5, arc=l, delta=10](Kappa,A,Lambda)
        \tkzLabelAngle[pos=0.8](Kappa,A,Lambda){$\omega$}

        % Γωνία ΣBΡ = θ (γύρω από το B)
        \tkzMarkAngle[fill=gray!30, opacity=0.8, size=0.5, arc=l, delta=10](Sigma,B,Rho)
        \tkzLabelAngle[pos=0.8](Sigma,B,Rho){$\theta$}

        % Γωνία MΓN = φ (γύρω από το Γ)
        \tkzMarkAngle[fill=gray!30, opacity=0.8, size=0.5, arc=l, delta=10](Mu,Gamma,Nu)
        \tkzLabelAngle[pos=0.8](Mu,Gamma,Nu){$\varphi$}

        % Σχεδίαση των καμπυλόγραμμων τμημάτων του ιμάντα (τόξα)
        % \tkzDrawArc(Center, StartPoint)(EndPoint) σχεδιάζει αντιωρολογιακά.
        % Τόξο από K προς Λ γύρω από A (από κάτω-αριστερά προς πάνω-αριστερά, αντιωρολογιακά)
        \tkzDrawArc[line width=2pt, color=maincolor](A,Kappa)(Lambda)
        % Τόξο από M προς N γύρω από Γ (από πάνω-δεξιά προς κάτω-δεξιά, αντιωρολογιακά)
        \tkzDrawArc[line width=2pt, color=maincolor](Gamma,Mu)(Nu)
        % Τόξο από P προς Σ γύρω από B (από κάτω-δεξιά προς κάτω-αριστερά, αντιωρολογιακά)
        \tkzDrawArc[line width=2pt, color=maincolor](B,Rho)(Sigma)

    \end{tikzpicture}
\end{wrapfigure}
Στο παρακάτω σχήμα τρεις κυκλικοί τροχοί με ίσες ακτίνες μήκους R, έχουν τα κέντρα τους στις κορυφές τριγώνου ΑΒΓ με πλευρές $\alpha$, $\beta$ και $\gamma$. Ένας τεντωμένος ιμάντας μήκους L συνδέει τους τρεις ίσους τροχούς όπως στο σχήμα και εφάπτεται σε αυτούς στα σημεία Κ, Λ, Μ, Ν, Ρ, Σ.

\begin{alist}
    \item Να αποδείξετε ότι:
    \begin{rlist}
        \item Το τετράπλευρο ΑΛΜΓ είναι ορθογώνιο. \monades{4}
        \item Η κυρτή γωνία $\mathrm{K\hat{A}\Lambda}$ και η γωνία $\hat{\mathrm{A}}$ του τριγώνου ΑΒΓ είναι παραπληρωματικές. \monades{4}
    \end{rlist}
    \item Αν $\mathrm{K\hat{A}\Lambda} = \hat{\omega}$, $\mathrm{\Sigma\hat{B}P} = \hat{\theta}$, $\mathrm{M\hat{\Gamma}N} = \hat{\varphi}$, να αποδείξετε ότι $\hat{\omega} + \hat{\theta} + \hat{\varphi} = 360^\circ$. \monades{8}
    \item Να αποδείξετε ότι το μήκος του ιμάντα L είναι $\mathrm{L = 2(\tau + \pi R)}$ όπου $\tau$ είναι η ημιπερίμετρος του τριγώνου ΑΒΓ. \monades{9}
\end{alist}